% Project-composed ordinary-world reading.  Every passage is class P.

\begin{courseunit}{M01}
\chapter{Reader: the brown fox}

\textbf{Text class: P --- project-composed instructional Latin.} No sentence
below is a quotation or an imitation attributed to an ancient author.

\section{Read before analyzing}

\Latin{Vulpes fusca est. Vulpes parva per silvam ambulat. Silva obscura est,
sed sol clarus est. Vulpes aquam videt. Aqua frigida est. Vulpes aquam bibit;
deinde sub arbore sedet.}

\textbf{First gloss.} \Latin{vulpes}, fox; \Latin{fuscus, -a, -um}, brown;
\Latin{parvus}, small; \Latin{silva}, wood; \Latin{obscurus}, dark;
\Latin{sol}, sun; \Latin{clarus}, bright; \Latin{aqua}, water;
\Latin{frigidus}, cold; \Latin{arbor}, tree; \Latin{ambulat}, walks;
\Latin{videt}, sees; \Latin{bibit}, drinks; \Latin{sedet}, sits;
\Latin{deinde}, then; \Latin{per} takes an accusative, and \Latin{sub} here
takes an ablative.

\section{What the endings are already doing}

In \Latin{vulpes fusca est}, \Latin{vulpes} names what the statement is about,
so it is nominative; \Latin{fusca} agrees with that feminine subject.  In
\Latin{vulpes aquam videt}, the fox remains the subject, while \Latin{aquam}
is the thing directly seen and is therefore accusative.  An ending is evidence
for a function, not a translation by itself.

Mark every finite verb, draw one line to its subject, and box each word that
agrees with a noun.  Then change \Latin{aquam} to \Latin{silvam}: what does the
fox now see?  Change \Latin{fusca} to \Latin{fuscae} only after making the
subject plural: \Latin{vulpes fuscae sunt}.

\section{Memory and composition}

Cover the passage and recover these pairs: fox/brown, wood/dark, sun/bright,
water/cold.  Write four new copular sentences and two subject--verb--object
sentences.  Keep an agreement ledger: noun, gender, number, case, agreeing
word, and the evidence for each decision.

\ifsolutions
\section{Models and points to check}

\textbf{Project translation.} The fox is brown. A small fox walks through the
wood. The wood is dark, but the sun is bright. The fox sees water. The water is
cold. The fox drinks the water; then it sits beneath a tree.

Possible fresh sentences are \Latin{silva magna est}; \Latin{aqua clara est};
\Latin{puella vulpem videt}; and \Latin{vulpes puellam videt}.  In the last two,
the endings, not position alone, distinguish the seer from the thing seen.
\else
\section{Before continuing}

You should reasonably be able to find the subject, finite verb, complement,
and direct object in fresh clauses of this size; explain basic agreement; and
compose six accurate ordinary-world sentences without copying the models.
\fi
\end{courseunit}

\begin{courseunit}{M02}
\chapter{Reader: house, garden, and field}

\textbf{Text class: P --- project-composed instructional Latin.}

\section{Connected prose}

\Latin{Iulia in villa habitat. Villa hortum parvum habet. In horto rosae et
lilia crescunt. Marcus, filius Iuliae, cum amico ad agrum ambulat. Agricola
servum vocat et pueris mala dat. Pueri dona portant; servus autem aquam ex
puteo portat. Sub vesperum mater filium et amicum domum vocat.}

\textbf{Cumulative gloss.} \Latin{villa}, country house; \Latin{hortus}, garden;
\Latin{rosa}, rose; \Latin{lilium}, lily; \Latin{ager, agri}, field;
\Latin{agricola}, farmer; \Latin{servus}, servant; \Latin{puer}, boy;
\Latin{malum}, apple; \Latin{donum}, gift; \Latin{puteus}, well;
\Latin{vesper}, evening; \Latin{habitat}, lives; \Latin{habet}, has;
\Latin{crescunt}, grow; \Latin{vocat}, calls; \Latin{dat}, gives;
\Latin{portat}, carries.  Here \Latin{cum} and \Latin{ex} take the ablative,
while \Latin{ad} takes the accusative.

\section{Recover three noun systems}

Sort the nouns into first, second masculine, and second neuter patterns.  Then
parse \Latin{hortum}, \Latin{rosae}, \Latin{agrum}, \Latin{amico},
\Latin{pueris}, \Latin{mala}, and \Latin{puteo}.  The neuter reminder belongs
in memory: nominative and accusative match; their plural ends in \Latin{-a}.

Rebuild these transformations without looking back:

\begin{enumerate}
\item one apple becomes many apples;
\item a boy gives a gift to the farmer;
\item the servants carry water out of the wells;
\item Julia walks toward the field with her son.
\end{enumerate}

\section{Retell from a map}

Draw four places---house, garden, field, well---and label them in Latin.  Retell
the paragraph from the map, first in the present and then with one new person.
The map should prompt meaning; it should not reproduce the original wording.

\ifsolutions
\section{Models and points to check}

\textbf{Project translation.} Julia lives in a country house. The house has a
small garden. Roses and lilies grow in the garden. Marcus, Julia's son, walks
to the field with a friend. The farmer calls a servant and gives apples to the
boys. The boys carry gifts, while the servant carries water from the well. At
evening the mother calls her son and his friend home.

Transformation models: \Latin{mala}; \Latin{puer agricolae donum dat};
\Latin{servi aquam ex puteis portant}; \Latin{Iulia cum filio ad agrum
ambulat}.  Other word orders are possible when the endings remain coherent.
\else
\section{Before continuing}

You should reasonably be able to recover the first- and second-declension
paradigms, recognize a case from both ending and syntax, and retell this scene
in a connected paragraph with dependable number and case.
\fi
\end{courseunit}

\begin{courseunit}{M03}
\chapter{Reader: road, town, and weather}

\textbf{Text class: P --- project-composed instructional Latin.}

\section{A day of travel}

\Latin{Prima luce viator e domo procedit. Via longa inter montes viatorem ad
urbem ducit. Viator sarcinam gravem manu portat, sed canis fidelis iuxta eum currit.
Meridie nubes atrae caelum tegunt; ventus fortis arbores movet, et imber viam
implet. Viator spem non amittit. Post breve tempus nubes abeunt, facies caeli
clarescit, et viator portam urbis conspicit.}

\textbf{Cumulative gloss.} \Latin{lux, lucis}, light; \Latin{viator}, traveler;
\Latin{domus}, house; \Latin{via}, road; \Latin{mons, montis}, mountain;
\Latin{urbs, urbis}, town; \Latin{sarcina}, pack; \Latin{manus}, hand;
\Latin{canis}, dog; \Latin{meridies}, midday; \Latin{nubes}, cloud;
\Latin{caelum}, sky; \Latin{ventus}, wind; \Latin{imber}, rain;
\Latin{spes}, hope; \Latin{tempus}, time; \Latin{facies}, face or appearance;
\Latin{porta}, gate; \Latin{procedit}, sets out; \Latin{ducit}, leads;
\Latin{currit}, runs; \Latin{tegunt}, cover; \Latin{movet}, moves;
\Latin{implet}, fills; \Latin{amittit}, loses; \Latin{abeunt}, depart;
\Latin{clarescit}, clears; \Latin{conspicit}, catches sight of.

\section{Complete the noun map}

Find at least one third-, fourth-, and fifth-declension noun.  Decline
\Latin{mons}, \Latin{manus}, and \Latin{spes} in the singular; then give the
forms demanded by these phrases: ``between the mountains,'' ``with both
hands,'' ``the hopes of the travelers,'' and ``the faces of the townspeople.''

Agreement can cross unlike-looking endings: \Latin{sarcinam gravem},
\Latin{canis fidelis}, \Latin{nubes atrae}, \Latin{ventus fortis}.  For each
pair, state gender, number, case, and function.  Do not demand matching letters
where Latin demands matching grammatical features.

\section{Compose weather that changes action}

Write eight to ten connected sentences in which weather interrupts a journey
or a day's work.  Include one noun from every declension, four adjective--noun
pairs, and at least two prepositional phrases.  Underline the exact ending that
supports each parse you would want to defend later.

\ifsolutions
\section{Models and points to check}

\textbf{Project translation.} At first light a traveler sets out from home. A
long road between mountains leads the traveler to the town. The traveler carries a heavy
pack in his hand, but a faithful dog runs beside him. At midday dark clouds
cover the sky; a strong wind moves the trees, and rain fills the road. The
traveler does not lose hope. After a short time the clouds depart, the face of
the sky clears, and the traveler catches sight of the city gate.

Requested forms include \Latin{inter montes}, \Latin{ambabus manibus},
\Latin{spes viatorum}, and \Latin{facies oppidanorum}.  A sound composition
may differ completely from the passage while preserving the named controls.
\else
\section{Before continuing}

You should reasonably be able to recognize all five declensions in connected
prose, maintain adjective agreement across dissimilar endings, and compose a
short narrative in which nouns retain clear functions through several clauses.
\fi
\end{courseunit}

\begin{courseunit}{M04}
\chapter{Reader: what the town can do}

\textbf{Text class: P --- project-composed instructional Latin.}

\section{Present action and ability}

\Latin{Urbs parva est, sed multa habet. In foro mercatores panes et fructus
vendunt. Faber in officina laborat et mensas facit. Scriba epistulas scribit.
Medicus aegrotos visitat. Cives aquam ex fonte ducere possunt, et nautae merces
per flumen portare possunt. Hodie pons clausus est; ideo plaustra flumen
transire non possunt. Cives tamen laborant, consilium capiunt, atque alteram
viam parant.}

\textbf{Cumulative gloss.} \Latin{forum}, market; \Latin{mercator}, merchant;
\Latin{panis}, bread; \Latin{fructus}, fruit; \Latin{faber}, craftsperson;
\Latin{officina}, workshop; \Latin{mensa}, table; \Latin{scriba}, scribe;
\Latin{epistula}, letter; \Latin{medicus}, physician; \Latin{aegrotus}, sick
person; \Latin{fons, fontis}, spring; \Latin{nauta}, sailor;
\Latin{merx, mercis}, merchandise; \Latin{flumen}, river; \Latin{pons}, bridge;
\Latin{plaustrum}, wagon; \Latin{consilium capere}, form a plan;
\Latin{vendunt}, sell; \Latin{facit}, makes; \Latin{visitat}, visits;
\Latin{ducere}, lead; \Latin{transire}, cross; \Latin{parant}, prepare.

\section{Finite verb, infinitive, and complement}

For every finite verb, give person, number, and lexical form.  In each phrase
with \Latin{possunt}, keep the infinitive attached to the idea of ability:
\Latin{cives possunt} is incomplete in this context; \Latin{cives aquam ducere
possunt} states what they can do.  Contrast \Latin{pons clausus est}, where an
adjective or participle is predicated of the bridge.

Turn the paragraph into a ledger with four columns: person, ordinary work,
finite verb, object or complement.  Then write a six-sentence account of what
people in another place are, do, and can or cannot do.

\ifsolutions
\section{Models and points to check}

\textbf{Project translation.} The town is small, but it has many things. In
the market merchants sell bread and fruit. A craftsperson works in a workshop
and makes tables. A scribe writes letters. A physician visits the sick. The
citizens can lead water from a spring, and sailors can carry merchandise along
the river. Today the bridge is closed, so wagons cannot cross the river. The
citizens nevertheless work, form a plan, and prepare another road.

One model extension: \Latin{In vico agricolae habitant. Frumentum colunt et
ad urbem portare possunt. Via autem lutosa est. Plaustra celeriter procedere
non possunt. Agricolae patientes sunt et equos ducunt. Ante noctem ad forum
veniunt.}
\else
\section{Before continuing}

You should reasonably be able to produce the active present system and forms
of \Latin{sum} and \Latin{possum}, keep infinitives distinct from finite verbs,
and sustain a present-tense description without slipping among persons.
\fi
\end{courseunit}

\begin{courseunit}{M05}
\chapter{Reader: the lost letter}

\textbf{Text class: P --- project-composed instructional Latin.}

\section{A narrative across time and voice}

\Latin{Heri scriba epistulam gravem magistratui misit. Tabellarius ante lucem
profectus est, sed in via tempestate oppressus est. Equus territus substitit;
epistula e sacco lapsa est et a vento sub saxum tracta est. Tabellarius diu
quaerebatur a civibus, qui de mora mirabantur. Tandem pastor advenit. Ille
epistulam invenerat et tabellarium ad villam duxit. Hodie epistula magistratui
reddita est; cras responsum scribetur.}

\textbf{Cumulative gloss.} \Latin{magistratus}, magistrate; \Latin{tabellarius},
courier; \Latin{tempestas}, storm; \Latin{equus}, horse; \Latin{saccus}, bag;
\Latin{ventus}, wind; \Latin{saxum}, rock; \Latin{mora}, delay;
\Latin{pastor}, shepherd; \Latin{proficiscor, proficisci, profectus sum}, set
out; \Latin{labor, labi, lapsus sum}, slip or fall; \Latin{miror, mirari,
miratus sum}, wonder; \Latin{misit}, sent; \Latin{oppressus est}, was overtaken;
\Latin{substitit}, stopped; \Latin{tracta est}, was dragged;
\Latin{quaerebatur}, was being sought; \Latin{invenerat}, had found;
\Latin{reddita est}, was returned; \Latin{scribetur}, will be written.

\section{Build the time line}

Place every verb under before-the-story, main past, ongoing past, present
result, or future.  Then distinguish passive forms from deponents: the courier
\emph{set out} in \Latin{profectus est}; he was not ``set out by someone.''  The
letter \Latin{lapsa est} likewise fell; but \Latin{tracta est} is genuinely
passive and names its agent with \Latin{a vento}.

Retell the event from the shepherd's viewpoint.  Use at least one imperfect,
one perfect, one pluperfect, one passive, and one deponent.  Then reduce the
story to three Latin sentences without losing the causal order.

\ifsolutions
\section{Models and points to check}

\textbf{Project translation.} Yesterday a scribe sent an important letter to
the magistrate. The courier set out before dawn, but was overtaken by a storm
on the road. His frightened horse stopped; the letter fell from the bag and
was dragged by the wind beneath a rock. The courier was sought for a long time
by citizens, who wondered about the delay. At last a shepherd arrived. He had
found the letter and led the courier to a farmhouse. Today the letter has been
returned to the magistrate; tomorrow a reply will be written.

A compressed model: \Latin{Tabellarius profectus est, sed tempestas eum
oppressit et epistula lapsa est. Pastor, qui epistulam invenerat, tabellarium
adiuvit. Epistula reddita est et responsum scribetur.}
\else
\section{Before continuing}

You should reasonably be able to maintain a coherent time line through the
indicative system, distinguish passive from deponent form and meaning, state
agency accurately, and retell a narrative rather than merely replace words.
\fi
\end{courseunit}

\begin{courseunit}{M06}
\chapter{Reader: from an ordinary message to a formal one}

\textbf{Text class: P --- project-composed instructional Latin, except for the
short received Missal dialogue explicitly labeled R below.}

\section{A message with shifting reference}

\Latin{Lucia fratrem suum ad portam invenit. Ei parvum volumen tradit et dicit:
``Hoc Marco da; ille rem exspectat.'' Frater eam interrogat: ``Ubi eum
inveniam?'' Lucia respondet: ``Cum amicis suis apud curiam laborat. Ad eos
celeriter ambula, sed volumen apud te clausum tene.'' Frater igitur per forum
ad curiam procedit. Ibi Marcum salutat eique volumen tradit. Marcus id aperit,
verba legit, atque Luciae per fratrem gratias agit.}

\textbf{Cumulative gloss.} \Latin{frater}, brother; \Latin{porta}, gate;
\Latin{volumen}, roll or document; \Latin{curia}, council house;
\Latin{trado}, hand over; \Latin{exspecto}, await; \Latin{interrogo}, ask;
\Latin{inveniam}, I may find; \Latin{teneo}, keep; \Latin{saluto}, greet;
\Latin{aperio}, open; \Latin{lego}, read; \Latin{gratias ago}, give thanks.
Pronoun ledger: \Latin{ei/eique} refers to Marcus or the brother as context
states; \Latin{eam} to Lucia; \Latin{eum/ille} to Marcus; \Latin{eos} to
Marcus and his companions; \Latin{id/hoc} to the document.

\section{One reader method}

Use the same five passes whenever prose becomes unfamiliar:

\begin{enumerate}
\item mark finite verbs and clause boundaries;
\item identify subjects, complements, and governed cases;
\item resolve pronouns from agreement, syntax, and sense;
\item notice tense, voice, connectors, and marked order;
\item make a direct first translation, then revise only where the Latin gives
evidence.
\end{enumerate}

Apply the passes to the dialogue, then to this received liturgical exchange
from the Ordinary of the 1962 Missal (R): \Latin{Dominus vobiscum. Et cum
spiritu tuo.}  The vocabulary is different, but the grammatical questions are
the same: implied verb, predicative or associative phrase, person addressed,
and the antecedent supplied by the ritual exchange.

\section{Reconstruction and transfer}

Close the book and reconstruct the message from these cues: Lucia---brother---
document---Marcus---council house---thanks.  Then compose a different ten-line
exchange using at least four pronouns and three prepositions.  Annotate every
pronoun with its antecedent and every preposition with its governed case.

\ifsolutions
\section{Models and points to check}

\textbf{Project translation.} Lucia finds her brother at the gate. She hands
him a small document and says, ``Give this to Marcus; he is waiting for the
matter.'' Her brother asks her, ``Where shall I find him?'' Lucia replies, ``He
is working with his friends at the council house. Walk quickly to them, but
keep the document closed with you.'' Her brother therefore goes through the
market to the council house. There he greets Marcus and hands him the document.
Marcus opens it, reads the words, and thanks Lucia through her brother.

\textbf{English rendering status for R: unavailable.}
The Latin exchange is elliptical.  That is an
interpretive observation, not permission to insert any convenient verb into
an unfamiliar clause.
\else
\section{Before continuing}

You should reasonably be able to resolve reference across a paragraph, explain
prepositional government, use the five-pass method on unfamiliar neutral and
liturgical clauses, reconstruct the sense from sparse cues, and compose a
coherent exchange whose pronouns remain unambiguous.
\fi
\end{courseunit}
